(1) 记 \(k_c=\pi/a\)。导体壁 \(x=0,a\) 上切向 \(E_y=0\) 已由正弦因子满足；在 \(y=0,b\) 上 \(E_y\) 为法向分量，亦不违反切向条件。以 \(e^{-i\omega t}\) 约定，Faraday 定律是 \(\nabla\times\mathbf E=i\omega\mu_0\mathbf H\)。逐分量求导得到
\[
H_x=-\frac{\beta}{\omega\mu_0}E_y,\qquad
H_z=-\frac{i k_c E_0}{\omega\mu_0}\cos(k_cx)e^{i(\beta z-\omega t)},\qquad H_y=E_x=E_z=0.
\]
再将它们代入 \(\nabla\times\mathbf H=-i\omega\epsilon_0\mathbf E\)，得到 \(\beta^2+k_c^2=\omega^2\epsilon_0\mu_0=\omega^2/c^2\)。一般矩形波导 TE 模有 \(k_c^2=(m\pi/a)^2+(n\pi/b)^2\)，不允许 \(m=n=0\) 的非平凡 TE 模；因 \(a>b\)，TE10 的 \(k_c=\pi/a\) 最小。电场、磁场相位符号与上述约定严格对应。

(2) 传播要求 \(\beta>0\) 为实数，故
\[
\omega_c=\frac{\pi c}{a},\qquad
\beta=\frac1c\sqrt{\omega^2-\omega_c^2},\quad
v_\phi=\frac{\omega}{\beta},\quad
v_g=\left(\frac{d\beta}{d\omega}\right)^{-1}=\frac{c^2\beta}{\omega}.
\]
由此 \(v_\phi v_g=c^2\)；相速度可以大于 \(c\)，但群速度 \(v_g=c\sqrt{1-(\omega_c/\omega)^2}<c\)，不违反因果律。TE 波阻抗定义为横向电场与横向磁场振幅比，得 \(Z_{\rm TE}=\omega\mu_0/\beta\)。当 \(\omega\downarrow\omega_c\)，\(v_g\to0,\ Z_{\rm TE}\to\infty\)；当 \(\omega\to\infty\)，\(v_\phi,v_g\to c,\ Z_{\rm TE}\to\sqrt{\mu_0/\epsilon_0}\)。

(3) 复场的时间平均能流 \(\langle\mathbf S\rangle=\frac12\operatorname{Re}(\mathbf E\times\mathbf H^*)\)。由于 \(H_x=-\beta E_y/(\omega\mu_0)\)，有
\[
\langle S_z\rangle=\frac{\beta E_0^2}{2\omega\mu_0}\sin^2(k_cx),\qquad
\boxed{P=\int_0^b\!\int_0^a\langle S_z\rangle\,dx\,dy
=\frac{ab\beta E_0^2}{4\omega\mu_0}}.
\]
单位长度的时间平均场能为
\[
U=\int_A\frac{\epsilon_0|\mathbf E|^2+\mu_0|\mathbf H|^2}{4}\,dA.
\]
利用 \(\int_0^a\sin^2(k_cx)dx=\int_0^a\cos^2(k_cx)dx=a/2\) 与 \(\beta^2+k_c^2=\omega^2\epsilon_0\mu_0\)，电能和磁能各为 \(\epsilon_0abE_0^2/8\)，故
\[
\boxed{U=\frac{\epsilon_0abE_0^2}{4}},\qquad
\boxed{\frac PU=\frac{\beta}{\omega\epsilon_0\mu_0}
=\frac{c^2\beta}{\omega}=v_g}.
\]
这里 \(P\) 的单位为瓦特，\(U\) 为焦耳每米，因此 \(P/U\) 具有速度量纲。

(4) 直接两次微分 \(\beta(\omega)=c^{-1}\sqrt{\omega^2-\omega_c^2}\)，得到
\[
\boxed{\beta''(\omega)=-\frac{\omega_c^2}{c(\omega^2-\omega_c^2)^{3/2}}<0}.
\]
在 \(\omega_0>\omega_c\) 附近展开 \(\beta=\beta_0+\beta_1\Omega+\tfrac12\beta_2\Omega^2\)，其中 \(\Omega=\omega-\omega_0,\ \beta_1=1/v_g,\ \beta_2=\beta''(\omega_0)\)。高斯积分给出随 \(z\) 传播的包络
\[
\mathcal E(z,t)\propto
\frac{\exp\!\left[-\frac{\sigma^2(t-\beta_1z)^2}{2(1-i\beta_2z\sigma^2)}\right]}
{\sqrt{1-i\beta_2z\sigma^2}},
\]
因此强度的均方根时间宽度为
\[
\boxed{\Delta t(z)=\frac{\sqrt{1+(\beta_2z\sigma^2)^2}}{\sqrt2\,\sigma}}.
\]
这要求 \(\sigma\ll\omega_0-\omega_c\)，以保证二阶近似及传播频谱的有效性。低于截止时令 \(\beta=i\kappa,\ \kappa=\sqrt{\omega_c^2-\omega^2}/c>0\)，半无限波导选择 \(e^{-\kappa z}\) 衰减解；\(H_x=-i\kappa E_y/(\omega\mu_0)\) 与 \(E_y\) 正交相位，故 \(\langle S_z\rangle=0\)。这并不否认有限长度倏逝区可与两端传播模耦合而发生隧穿透射。
